\documentclass[12pt]{article}
\usepackage[utf8]{inputenc}
\usepackage[T1]{fontenc}
\usepackage{hyperref}
\title{Soyaltepec Mixtec Language Resources for K-12 Curriculum Development}
\author{Compiled by Claude Sonnet 5 (claude-sonnet-5) for Andr\'as Kornai}
\date{}
\begin{document}
\maketitle
\section{Introduction}
Soyaltepec Mixtec (ISO 639-3 vmq, Glottocode soya1236) is an Otomanguean language of Oaxaca, Mexico, spoken in San Bartolo Soyaltepec and Guadalupe Gabilera. Wikitongues describes it as unwritten and rates it EGIDS 8a (moribund). ASJP lists about 220 speakers. The search found no school materials, grammar or dictionary.
\section{Primary Educational Resources}
None found.
\section{Linguistic and Reference Materials}
\begin{itemize}
\item Wikipedia entry for ISO 639 vmq: \url{https://wikipedia.com/wiki/ISO_639:vmq}
\item ASJP wordlist, San Bartolo Soyaltepec Mixtec; it cites Josserand's 1983 dissertation \emph{Mixtec Dialect History}: \url{https://asjp.clld.org/languages/SAN_BARTOLO_SOYALTEPEC_MIXTEC}
\item AILLA (UT Austin), 1975 ``List of 500'' wordlist with audio, San Bartolo Soyaltepec (possibly vmq): \url{https://islandora-ailla.lib.utexas.edu/islandora/object/ailla%3A129083}
\item Wikitongues language page: \url{https://wikitongues.org/languages/vmq/}
\end{itemize}
\section{Religious and Community Resources}
None found in the search results.
\section{Digital and Audio Resources}
\begin{itemize}
\item Global Recordings Network, Mixteco de Soyaltepec: \url{https://globalrecordings.net/ja/language/vmq}
\end{itemize}
\section{Curriculum Development Context}
The language is moribund and unwritten, and the search found almost no existing materials. Any curriculum would begin with documentation: an orthography, recording elders and building a lexicon from the archived wordlist.
\end{document}